\chapter*{Abstract}
This report documents the practical exploitation and runtime validation of logic flaws in the Linux eBPF verifier on LTS kernel 6.8. Building on a prior ISO-IEC TS 17961-2013~\cite{ISO17961} rule-based assessment conducted by a previous research team, we implement \textbf{proof-of-concept (PoC) exploits} and collect \textbf{multi-stage evidence} (bytecode, verifier logs, translated code, and runtime traces) to bridge the gap between theoretical verifier bypass analysis and actual runtime impact. \textbf{Five vulnerabilities} are analyzed in depth through systematic exploitation: three \textbf{ABI/argument-mismatch cases} (5.06a/5.06b/5.06d) representing function pointer signature violations, one unprototyped function pointer case with \textbf{stack out-of-bounds write} (5.39) enabling \textbf{eBPF-local memory corruption} (sandbox-limited, not kernel-level), and one \textbf{tainted-length packet buffer overflow} (5.46b) corrupting packet-local memory (bounds-checking gap, not kernel compromise). The work demonstrates which verifier failures manifest merely as observable runtime behaviors (triggers) versus genuine memory-corruption primitives, presents a \textbf{reproducible evidence-chain methodology} for eBPF vulnerability assessment, and documents the design and deployment of an isolated external-access testing environment enabling real-world-like exploitation scenarios from the host machine to the target VM.
\clearpage
